\color{unimain}{\subsection{Газовая защита трансформатора (ГЗ Т)}}
\color{black}

\begin{enumerate}[label=\arabic{section}.\arabic{subsection}.\arabic*,  align=left, labelsep=4pt, leftmargin=0pt, itemindent=57pt]

\item
Газовая защита трансформатора предназначена для защиты от повреждений внутри кожуха,
сопровождающихся выделением газа, и понижением уровня масла.  

\begin{figure}[H]
\centering
\includegraphics[width=0.88\textwidth,height=0.88\textheight,keepaspectratio]{\fbpath/07.01 РЗТ/011.1071(1072) ГЗ Т/img/GZT_sign.pdf}
\captionof{figure}{Функциональная схема ГЗ Т сигнальная ступень}\label{pict:GZT_sign}
\end{figure}

\begin{figure}[H]
\centering
\includegraphics[width=0.88\textwidth,height=0.88\textheight,keepaspectratio]{\fbpath/07.01 РЗТ/011.1071(1072) ГЗ Т/img/GZT_otkl.pdf}
\captionof{figure}{Функциональная схема ГЗ Т отключающая ступень}\label{pict:GZT_otkl}
\end{figure}


\item
Предусмотрена возможность приема сигналов срабатывания сигнального контакта ГЗ Т, отключающего
контакта ГЗ Т и низкой изоляции цепи отключающего контакта ГЗ Т, по разным цепям от двух источников ПДС1
и ПДС2. Для приема сигналов от ПДС1 предусмотрены входные сигналы \textit{«Ср. сигн контакта ГЗ Т 1 цепь»},
\textit{«Ср. откл контакта ГЗ Т 1 цепь»} и \textit{«Низкая изол. ГЗ Т откл 1 цепь»}, а для приема сигналов от ПДС2
предусмотрены входные сигналы \textit{«Ср. сигн контакта ГЗ Т 2 цепь»}, \textit{«Ср. откл контакта ГЗ Т 2 цепь»} и
\textit{«Низкая изол. ГЗ Т откл 2 цепь»}. Если сигналы срабатывания ГЗ Т заводятся непосредственно со шкафа ШОМ
силового трансформатора конфигурируются только входные сигналы \textit{«Ср. сигн контакта ГЗ Т 1 цепь»}, \textit{«Ср.
откл контакта ГЗ Т 1 цепь»} и \textit{«Низкая изол. ГЗ Т откл 1 цепь»}, а действие ГЗ Т по второй цепи сигнального
и отключающего контакта выводиться из работы с помощью программных переключателей \textbf{«ГЗ Т сигн 2 цепь»}
и \textbf{«ГЗ Т откл 2 цепь»}.

\item
Ввод в работу ГЗ Т осуществляется с помощью программных переключателей \textbf{«ГЗ Т сигн 1 цепь»}, \textbf{«ГЗ
Т сигн 2 цепь»}, \textbf{«ГЗ Т откл 1 цепь»} и \textbf{«ГЗ Т откл 2 цепь»}. Оперативный перевод ГЗ Т на сигнал производится
с помощью команд управления \textit{«Перевод ГЗ Т сигн на сигнал»} и \textit{«Перевод ГЗ Т откл на сигнал»}.
Предусматривается возможность действия срабатывания сигнального контакта ГЗ Т на отключение
трансформатора с помощью программного переключателя \textbf{«Откл от ГЗ Т сигн»}.

\item
Оперативный ввод/вывод ГЗ Т сигнальной ступени производится с помощью БУ (блока управления
функциями) \textit{«Вывод ГЗ Т сигн»}, а оперативный ввод/вывод ГЗ Т отключающий ступени производится с
помощью БУ (блока управления функциями) \textit{«Вывод ГЗ Т откл»}.

\item
Непрерывный контроль изоляции цепи отключающего контакта ГЗ Т осуществляется с помощью реле
контроля изоляции (РКИ). В случае снижения сопротивления изоляции РКИ подает сигнал \textit{«Низкая изол. ГЗ Т
откл 1 цепь»} (\textit{«Низкая изол. ГЗ Т откл 2 цепь»}) для блокировки срабатывания ГЗ Т от отключающего контакта.
Время срабатывания блокировки от сигналов РКИ задается с помощью уставки \textbf{«Тбл ГЗ Т откл»}. Ввод
блокировки ГЗ Т от РКИ осуществляется с программного переключателя \textbf{«Блок ГЗ Т откл»}. Для сброса
блокировки от РКИ, используется команда управления \textit{«Возврат ГЗ Т откл»}.


\small
\setlength{\extrarowheight}{0.05cm} % высота ячеек в таблице
\begin{longtable}
{|>{\centering\arraybackslash}m{4.5cm}|>{\centering\arraybackslash}m{4.5cm}|>{\centering\arraybackslash}m{3cm}|>{\centering\arraybackslash}m{3.5cm}|}
\caption{Перечень уставочных параметров ГЗ Т\hfill\vspace{-0.5\baselineskip}}\label{gzt:tbl1} \\ 
\hhline{|-|-|-|-|}
\rowcolor{gray!30}
Наименование уставки & Диапазон & Шаг &  Значение по умолчанию \\ 
\hhline{|-|-|-|-|}
\endfirsthead

\caption*{\hspace{3pt}\emph{Продолжение таблицы \ref{gzt:tbl1}\hfill\vspace{-0.5\baselineskip}}} \\ % сделано по ГОСТ 2.105 п.6.8.7
\hhline{|-|-|-|-|}
\rowcolor{gray!30}
Наименование уставки & Диапазон & Шаг &  Значение по умолчанию \\  
\hhline{|-|-|-|-|}
\endhead
\hhline{|-|-|-|-|}

\endfoot
\endlastfoot

%%%%%%%%%%%%содержимое%%%%%%%%%%%%%%

\end{longtable}
\setlength{\extrarowheight}{0.0cm} % возврат параметра высота ячеек в таблице
\normalsize










































\end{enumerate} 